\section*{第 \ref{ch:b3:quantitative-geomorphology} 章——定量地貌学}

\begin{solution}{exo:b3:quantitative-geomorphology:1}
$a/\sqrt{\pi\kappa t} = 0.30 - 0.04 = 0.26$，故 $\kappa t = (2/0.26)^2/\pi =
\qty{18.8}{m^2}$，且 $t = 18.8/(5\times10^{-3}) = \qty{3.8}{ka}$，
结合 $\kappa$ 的取值范围，两个方向上的误差均在两倍以内。
\end{solution}

\begin{solution}{exo:b3:quantitative-geomorphology:2}
$E = \kappa\,|\nabla^2 z| = 4\times10^{-3}\times0.012 =
\qty{4.8e-5}{m/\text{年}}$，即 \qty{0.048}{mm/\text{年}}。地形起伏 $E L^2/(2\kappa) =
0.012\times80^2/2 = \qty{38.4}{m}$.
\end{solution}

\begin{solution}{exo:b3:quantitative-geomorphology:3}
$E = 2\times10^{-6}\times(25\times10^6)^{0.5}\times0.01 = 2\times10^{-6}
\times5000\times0.01 = \qty{1.0e-4}{m/\text{年}}$，即 \qty{0.1}{mm/\text{年}}。由于 $E$ 的单位为
\unit{m/\text{年}}，$A^m$ 的单位为 $\unit{m}^{2m}$，$S$ 无量纲，故 $K$ 的单位为
$\unit{m}^{1-2m}\,\unit{\text{年}^{-1}}$：当 $m = 0.5$ 时为 \unit{\text{年}^{-1}}。
\end{solution}

\begin{solution}{exo:b3:quantitative-geomorphology:4}
(a) 固定型：硬岩床在出露处使河流变陡，且台阶
随岩床保持不动。(b) 锚定型：只要断层仍在滑动，
河流在跨越断层处即被变陡。(c) 移动型：海岸的隆升
使所有河流的\omterm{def:b3:quantitative-geomorphology:knickpoint}{侵蚀基准面}同时下降，\omterm{def:b3:quantitative-geomorphology:knickpoint}{裂点}
正在向上游迁移，并且正如运动波定理对 $n
= 1$ 的预言那样，它们位于同一条等高线上。
\end{solution}

\begin{solution}{exo:b3:quantitative-geomorphology:5}
$\theta = m/n = 0.5$. $S = (U/K)^{1/2}A^{-1/2} = (10^{-4}/10^{-9})^{1/2}
\times(10^7)^{-1/2} = 316\times3.16\times10^{-4} = 0.10$。加倍 $U$
使 $k_s = (U/K)^{1/n}$ 乘以 $\sqrt2 = 1.41$；当 $n = 1$ 时则乘以 2：
$n$ 越大，陡度对隆升的响应越弱。
\end{solution}

\begin{solution}{exo:b3:quantitative-geomorphology:6}
以 $x$ 表示距出口的上游距离，则 $\mathrm{d}x = -\mathrm{d}d$，故 $\chi(d)
= \int_d^L (A_0/c\,d'^{h})^{\theta}\,\mathrm{d}d' = A_0^{\theta} c^{-\theta}
\int_d^L d'^{-h\theta}\,\mathrm{d}d' = A_0^{\theta}c^{-\theta}(L^p -
d^p)/p$。此处 $p = 0.19$，$A_0^{\theta}c^{-\theta} = 10^{2.7}/2^{0.45}$，
且 $\chi = \qty{2935}{m}$。
\end{solution}

\begin{solution}{exo:b3:quantitative-geomorphology:7}
$c = 2\times10^{-6}\times(10^8)^{0.5} = \qty{0.02}{m/\text{年}}$, \qty{20}{m/ka};
以此速度通过 \qty{5}{km} 需 \qty{250}{ka}。面积及与之相关的
波速沿上游减小，因此\omterm{def:b3:quantitative-geomorphology:knickpoint}{裂点}上移时减速。
\end{solution}

\begin{solution}{exo:b3:quantitative-geomorphology:8}
$\bar e = 0.40\times600 = \qty{240}{m}$；山顶抬升 $0.80\times240 =
\qty{192}{m}$；平均海拔下降 $240 - 192 = \qty{48}{m}$。
\end{solution}

\begin{solution}{exo:b3:quantitative-geomorphology:9}
$h = h_0\ln(P_0/E)$：$0.5\ln 6 = \qty{0.90}{m}$，且 $0.5\ln1.5 =
\qty{0.20}{m}$。当侵蚀速率为 \qty{0.4}{mm/\text{年}} 时，侵蚀超过最快的
\omterm{def:b3:advanced-geochronology:cosmogenic}{产生率} $P_0$：土壤被剥离，坡面变为裸岩，其
降低速率受风化限制。
\end{solution}

\begin{solution}{exo:b3:quantitative-geomorphology:10}
\omterm{def:b3:quantitative-geomorphology:steady-profile}{稳态}下 $\mathrm{d}q/\mathrm{d}x = E$，且分水岭处 $q = 0$，故
$q = E x$。于是 $\kappa S = E x\,(1 - S^2/S_c^2)$，这是一个二次方程，其正根为
$S = [-\kappa + \sqrt{\kappa^2 + 4E^2x^2/S_c^2}\,]/(2Ex/S_c^2)$。在
20、50、100 和 \qty{200}{m} 处分别为 0.36、0.68、0.89 和 1.03，而线性扩散（$S = Ex/\kappa$）给出的对应值为 0.4、1.0、
2.0 和 4.0。山丘具有圆润的
山脊，随后是坡度趋近于 $S_c$ 的近似平直坡面，
其地形起伏远小于抛物线形的剖面。
\end{solution}

\begin{solution}{exo:b3:quantitative-geomorphology:11}
由运动波定理，变化抵达坐标为
$\chi$ 所需时间为 $\chi/(KA_0^m)$；河网完全响应，是指变化已
到达最后一个河道源头，即 $\chi$ 最大者。$K A_0^m = 10^{-6}\times
1000 = \qty{e-3}{m/\text{年}}$：主干需 \qty{5.46}{Ma} 响应，河网
需 \qty{7.24}{Ma}，与预期的约一百万年量级相符。
\end{solution}

\begin{solution}{exo:b3:quantitative-geomorphology:12}
在汇合点以上，支流的\omterm{def:b3:quantitative-geomorphology:steady-profile}{稳态}坡度为 $S = (U/2K)A^{-\theta}$，
故 $\mathrm{d}z/\mathrm{d}\chi = (U/2K)A_0^{-\theta}$，为主干的一半。
支流\omterm{def:b3:quantitative-geomorphology:knickpoint}{裂点}到达汇合点时，主干\omterm{def:b3:quantitative-geomorphology:knickpoint}{裂点}位于 $\chi_j$；此后
再经过时间 $\Delta t$，主干\omterm{def:b3:quantitative-geomorphology:knickpoint}{裂点}移至 $\chi_k$，且 $\chi_k -
\chi_j = K A_0^m\,\Delta t$，而支流\omterm{def:b3:quantitative-geomorphology:knickpoint}{裂点}的波速为 $2KA^m$，
则已移至 $\chi$，且 $\chi - \chi_j = 2KA_0^m\,\Delta t = 2(\chi_k -
\chi_j)$.
\end{solution}

\begin{solution}{pb:b3:quantitative-geomorphology:1}
\textbf{1.} $4.5\times12000^{1.7} = \qty{3.87e7}{m^2}$, \qty{38.7}{km^2};
$4.5\times300^{1.7} = \qty{7.3e4}{m^2}$, \qty{0.073}{km^2}.

\textbf{2.} $\theta \approx -\log(0.007855/0.06641)/\log(37.08/0.388) =
0.928/1.980 = 0.47$.

\textbf{3.} 是。十二个数据点的离散度约为 \qty{10}{\percent}，
$0.466$ 与 0.5 之间无显著差异。

\textbf{4.} $0.06641\times\sqrt{388\,000} = \qty{41.4}{m}$;
$0.007855\times\sqrt{3.708\times10^7} = \qty{47.8}{m}$。几何平均，
\qty{45.1}{m}，为 $k_s$ 的估计值，对应 $\theta = 0.5$。

\textbf{5.} 坡度是两个高程之差，二者各自
存在分米至米量级的误差，再除以较短的距离；在平缓河段上，
误差占高差的很大部分。$\chi$ 通过积分而非
微分得到。

\textbf{6.} $p = 1 - 1.7\times0.5 = 0.15$; $\chi = (1000/\sqrt{4.5})\,
(12000^{0.15} - 9000^{0.15})/0.15 = \qty{543.1}{m}$.

\textbf{7.} $(193.4 - 170)/543.1 = 0.0431$; $k_{sn} = 0.0431\times A_0^{0.5}
= \qty{43.1}{m}$.

\textbf{8.} χ 估算：它依据的是 \qty{3}{km} 长河段上的
高程，而非十二个带噪声的坡度。

\textbf{9.} $K = U/k_{sn} = 4.3\times10^{-5}/43.1 = \qty{1.00e-6}{\text{年}^{-1}}$.

\textbf{10.} $10^{-6}\times\sqrt{3.87\times10^7}\times0.00691 =
\qty{4.30e-5}{m/\text{年}}$，\qty{43}{m/Ma}：出口跟上了
隆升的速度。

\textbf{11.} 出口处 $10^{-6}\times6222 = \qty{6.2}{m/ka}$；
瀑布处，$d = \qty{7.40}{km}$，$A = \qty{1.70e7}{m^2}$，
速率为 \qty{4.1}{m/ka}。

\textbf{12.} $170 + 43\times(4.5)^{-0.5}(12000^{0.15} - 300^{0.15})/0.15 =
\qty{405}{m}$：地形起伏为 \qty{235}{m}。

\textbf{13.} 当 $n = 1$ 时，方程为运动学波：台阶两侧的
各点均以 $U$ 的速率抬升，因此台阶保持其高度，
即落差 \qty{25}{m}。

\textbf{14.} $\chi = (1000/\sqrt{4.5})(12000^{0.15} - 7400^{0.15})/0.15 =
\qty{899.4}{m}$.

\textbf{15.} $t = \chi/(KA_0^{0.5}) = 899.4/10^{-3} = \qty{899}{ka}$.

\textbf{16.} 当 $\chi = 899.4 + 300 = \qty{1199}{m}$ 时：
$d^{0.15} = 12000^{0.15} - 1199\times0.15\times\sqrt{4.5}/1000$, $d =
\qty{6.25}{km}$，位于出口上游 \qty{5.75}{km} 处。

\textbf{17.} 在汇合处（$d = \qty{9.5}{km}$）$\chi = \qty{442.8}{m}$；
瀑布位于支流的 $\chi$ 达到 \qty{899.4}{m} 处：
$s^{0.15} = 5000^{0.15} - (899.4 - 442.8)\times0.15\times\sqrt{1.5}/1000$,
$s = \qty{4.27}{km}$，位于汇合处上方 \qty{0.73}{km}。

\textbf{18.} 它们的汇合处位于 $\chi = 1954$ 与 \qty{999}{m}，均超过
\qty{899}{m}：\omterm{def:b3:quantitative-geomorphology:knickpoint}{裂点}尚未越过它们，因此该处的干流
仍保持原有高度。

\textbf{19.} $E = 5\times10^{-3}\times0.0086 = \qty{4.3e-5}{m/\text{年}}$,
\qty{43}{m/Ma}：瀑布上方的丘陵以隆升速率被侵蚀，处于
\omterm{def:b3:quantitative-geomorphology:steady-profile}{稳态}。

\textbf{20.} $150^2/(5\times10^{-3}) = \qty{4.5}{Ma}$，是
瀑布年龄的五倍。

\textbf{21.} 它们的坡脚随河道下降了 \qty{25}{m}，其
下部过陡，或许正在滑动，而坡顶仍
以原有速率被侵蚀。一道变陡的波沿每个坡面向上推进；经过几
百万年后，它们回到同一条抛物线，整体低了 \qty{25}{m}。

\textbf{22.} $t = \chi k_{sn}/(U A_0^{0.5})$ 与 $1/U$ 成正比：从
$899/1.2 = 749$ 到 $899/0.8 = \qty{1124}{ka}$。

\textbf{23.} 当 $n = 2$ 时，波速随坡度增大：陡峭的
坡面推进速度将快于周围河段，从而保持尖锐，但
其速度与由此读出的年龄将有所不同。

\textbf{24.} 瓦尔登河的\omterm{def:b3:quantitative-geomorphology:knickpoint}{侵蚀基准面}下降了
\textbf{\qty[text-series-to-math=true]{900}{ka}}前（按
计算为 \qty{899}{ka}）。
\end{solution}
